\paragraph{F.50 (Proved). Reconstruction over subfields of finitely generated cyclotomic extensions.}
Let $k$ be any subfield of a field finitely generated over
$\mathbb Q^{\mathrm{ab}}$. Then $k$ has pinned etale reconstruction:
for all finite-type $k$-schemes $X,Y$,
\[
 \operatorname{Hom}_k(X^{\mathrm{awn}},Y)
  \ \xrightarrow{\ \sim\ }\
 \pi_0\operatorname{Hom}^{\mathrm{pin}}_{B G_k}
             (X_{\mathrm{et}},Y_{\mathrm{et}}).
 \tag{F50}
\]
The relative commutation isomorphism is part of the morphism data.
This includes F.47 and also every subfield of a finitely generated
extension of an arbitrary cyclotomic extension of a number field.
We prove the additional arithmetic and quasi-rational steps that
extend F.47; its geometric and categorical deductions then apply
with all their hypotheses verified.

\begin{proof}
Choose $k\subset M$, where $M/\mathbb Q^{\mathrm{ab}}$ is finitely
generated, and choose an algebraic closure $\Omega$ of $M$. Let
$K\subset\Omega$ be the algebraic closure of $k$. In particular
$\overline{\mathbb Q}\subset K$. For every finite extension
$L/k$ inside $K$, the compositum $M_L=ML$ is finite over $M$,
hence again finitely generated over $\mathbb Q^{\mathrm{ab}}$.
Thus the field class in the statement is stable under every finite
extension needed for points, boundary components, or coefficient
descent.

\emph{1. Divisibility and the free multiplicative quotient.}
Let $E/\mathbb Q^{\mathrm{ab}}$ be finitely generated, and let
$E_0$ be the algebraic closure of $\mathbb Q^{\mathrm{ab}}$ in $E$.
This is a finite extension (Stacks, Lemma~9.26.11
\cite{StacksFGConstants}); hence $E_0=F_0\mathbb Q^{\mathrm{ab}}$
for a number field $F_0$. The extension $E/E_0$ is regular in
characteristic zero. Choose a proper normal geometrically integral
$E_0$-model $Z$ of $E$. Its divisor map gives an injection
\[
 E^*/E_0^*\lhook\joinrel\longrightarrow
       \operatorname{Div}(Z),                              \tag{F50.1}
\]
because a rational function with zero divisor on a proper normal
variety is a global unit, and its constant field is $E_0$.
The divisor group is free abelian, so its subgroup $E^*/E_0^*$ is
free abelian. By Tsujimura's Proposition~4.9 with parameter $f=1$,
$E_0^*/\mu_\infty$ is free abelian; the identification of its
divisible subgroup with $\mu_\infty$ is the Amoroso--Zannier
height argument in F.47. Consequently the exact sequence
\[
 0\longrightarrow E_0^*/\mu_\infty
 \longrightarrow E^*/\mu_\infty
 \longrightarrow E^*/E_0^*\longrightarrow0
\]
splits as a sequence of abelian groups, because its quotient is
free. Therefore $E^*/\mu_\infty$ is free abelian.

Also $D(E^*)=\bigcap_n(E^*)^n=\mu_\infty$. If $u$ is infinitely
divisible, its divisor is divisible by every integer, hence zero.
Every chosen $n$th root of $u$ has zero divisor as well, so $u$ and
all these roots lie in $E_0$. The height theorem applies there and
forces $u$ to be a root of unity. The converse holds since $E_0$
contains all roots of unity. In particular, for every finite $L/k$,
\[
               D(L^*)\subset\mu_\infty.                  \tag{F50.2}
\]
This follows by embedding $L$ and all its division witnesses into
$M_L$; no assertion that $L$ itself contains all roots of unity is
needed.

\emph{2. Abelian-variety Kummer faithfulness.}
For each finite $L/k$ and every abelian variety $A/L$,
\[
                    \bigcap_n nA(L)=0.                    \tag{F50.3}
\]
This is also the stability statement in Tsujimura,
Proposition~1.2(ii) \cite{Tsujimura2022}, combined with his
Theorem~3.1 over finite extensions of $\mathbb Q^{\mathrm{ab}}$
\cite{Tsujimura2020}. We give the specialization argument, including
the uniformity needed for all divisions.

For a finitely generated field $E/\mathbb Q^{\mathrm{ab}}$, an
abelian variety $A/E$, and a point $P\in\bigcap_n nA(E)$, spread
the abelian variety and this one point to an abelian scheme
$\mathcal A\to U$ and a section $P_U$, where $U$ is a normal
integral finite-type scheme over $\mathbb Q^{\mathrm{ab}}$ with
function field $E$. For every $n$, choose $Q_n\in A(E)$ with
$nQ_n=P$. The finite etale scheme $[n]^{-1}(P_U)\to U$ has the
generic section $Q_n$. Normality of $U$ extends it to a section on
the same $U$: the generic section selects a finite birational
component, which is isomorphic to the normal base. Thus all the
divisions extend without an infinite sequence of shrinkings.
At every closed point $x\in U$, the residue field is finite over
$\mathbb Q^{\mathrm{ab}}$, and $P_U(x)$ is divisible by all $n$.
Theorem~3.1 gives $P_U(x)=0$. Closed points are dense and the
abelian scheme is separated, so $P_U=0$ and $P=0$.
Apply this to $E=M_L$ and the base change of $A$; the injection
$A(L)\hookrightarrow A(M_L)$ proves (F50.3). The Picard and total
Tate-module conclusions of F.47 follow exactly as there, including
the treatment of Neron--Severi torsion.

\emph{3. Rational-source integrality without algebraic coefficients.}
The rational-source lemma of F.47 must now allow distinct punctures
$b_1,\ldots,b_r\in L$ which can be transcendental over $\mathbb Q$.
Put $E=\mathbb Q(b_1,\ldots,b_r)$, a finitely generated field.
There is an integral domain $A\subset E$, smooth and of finite
type over $\mathbb Z[1/N]$, with fraction field $E$, containing
all $b_i$ and all $(b_i-b_j)^{-1}$ for $i\ne j$. To construct it,
take a smooth affine model over $\mathbb Q$, shrink it to contain
these finitely many functions as regular functions and units, and
spread the finitely presented smooth algebra over $\mathbb Z[1/N]$;
Stacks, Lemma~10.138.14 \cite{StacksSmoothSpread} supplies this
last step. Flatness makes the model inject into its generic fibre,
so it remains a domain. Its smooth image in
$\operatorname{Spec}\mathbb Z[1/N]$ is an open set containing the
generic point, hence contains all but finitely many primes.

Choose different primes $p_i>r$ with nonempty fibres, and choose
a generic point $\mathfrak p_i$ of a component of the fibre over
$p_i$. The local ring $A_{\mathfrak p_i}$ is a discrete valuation
ring and $p_i$ is a uniformizer: the model is regular and flat and
the fibre is smooth, hence generically reduced. Let $v_i$ be this
valuation of $E$. Since the pairwise differences are units in
$A$, their residues at $\mathfrak p_i$ are nonzero. Among the
$p_i$ elements of $\mathbb F_{p_i}$ choose $c_i$ different from
all the at most $r$ residues of $b_j$ that lie in that prime field.
The Chinese remainder theorem in $A$ now supplies $x,y\in A$
such that
\[
 x\equiv b_i+p_i\pmod{p_i^2A},\qquad
 y\equiv c_i\pmod{p_iA}\quad(1\le i\le r).
\]
The ideals for different rational primes are comaximal. Therefore
\[
 v_i(x-b_j)=\delta_{ij},\qquad v_i(y-b_j)=0.                \tag{F50.4}
\]
In particular $x,y$ lie in the punctured affine line over $L$,
and the ratios $(x-b_j)/(y-b_j)$ are multiplicatively independent
modulo roots of unity. Independence persists in the overfield
$M_L$, since it is a statement about equalities of these same
elements.

Part~1 makes $M_L^*/\mu_\infty$ free abelian. The character-matrix
argument (F47.R4) therefore applies without alteration: finitely
many integer characters give a nonsingular integer matrix $N$ on
these independent ratios; two actual Kummer evaluations give
$N\lambda\in\mathbb Z^r$. The adjugate matrix and
$\mathbb Q\cap\widehat{\mathbb Z}=\mathbb Z$ force
$\lambda\in\mathbb Z^r$. A single further evaluation makes the
constant completed class an actual element of $L^*$.
Thus every compatible Kummer class on a punctured affine line
with actual specializations is an actual regular-unit class, in
the same generality as required for the norm calculation in F.47.
No valuation on the possibly much larger field $L$ was asserted
to be discrete.

\emph{4. The quasi-rational degree argument in the larger algebraic closure.}
Use the fixed subfield $B=\mathbb Q^{\mathrm{ab}}\subset K$.
All pointwise ambiguities furnished by (F50.2) lie in
$\mu_\infty\subset B^*$. The proof of the finite-exception lemma
in F.46 uses high-degree closed points after adjoining finitely
many coefficients. Here those coefficients generate a finitely
generated field $N/B$ inside $K$, which need not be algebraic.
We verify the required high-degree property directly.

Let $N_0=N\cap\overline{\mathbb Q}$. It is finite over $B$ by
the constants lemma, and $N/N_0$ is regular. Hence $N$ and
$\overline{\mathbb Q}$ are linearly disjoint over $N_0$; this is
the geometric irreducibility assertion in Stacks,
Lemma~10.47.13 \cite{StacksFGConstants}, using characteristic zero
for geometric reducedness. Thus, for algebraic numbers $a$,
\[
          [N(a):N]=[N_0(a):N_0].                           \tag{F50.5}
\]
The right side is unbounded. If
$\overline{\mathbb Q}/N_0$ had finite degree, its elements would
have uniformly bounded degree over $B$, and the Amoroso--Zannier
height bound would contradict the elements $2^{1/n}$.
Finite extensions of the infinite algebraic extension therefore
have unbounded degrees; characteristic zero supplies primitive
elements.

Descend the smooth curve and the finitely many rational functions
in a given translate argument to such an $N$. A finite map to
$\mathbb P^1_N$ and values $a$ as in (F50.5) give closed points
of arbitrarily large residue degree. A nonzero linear combination
of the fixed rational functions has a uniformly bounded zero
divisor, so cannot vanish at those points. The same choice avoids
any finite exceptional set: its points algebraic over $N$ have
bounded degrees, and its points transcendental over $N$ cannot
coincide with the chosen closed points. This is exactly the
degree argument used to obtain the rational identities, the two
covering outside-$B$ sets, and their infinite intersection in
F.46. Also $K/B$ is infinite dimensional, since it contains
$\overline{\mathbb Q}/B$, so the constant-factor case of that
proof applies. Consequently the lemma holds here without assuming
that $K/B$ is algebraic. We have verified the proof's actual
hypothesis, rather than using the algebraic-extension version of
Tsujimura's criterion beyond its stated range.

\emph{5. Completion of the reconstruction argument.}
The curve argument of F.47 now applies with its three inputs
verified for every finite $L/k$: the norm class on the punctured
line has integral residues by part~3; the integral divisor class
and the total Picard Tate module vanish by (F50.3); and the
pointwise quotient is a root of unity by (F50.2). Part~4 gives a
rational point function from all translates, including every
exceptional value. Galois equivariance descends that function,
and the rational-point Hochschild--Serre argument identifies its
entire compatible Kummer class.

The higher-dimensional argument of F.47 uses only characteristic
zero, these curve assertions after finite extensions, and
(F50.3). Smooth compactification and transverse curves make the
boundary residues integral; the Picard/Neron--Severi argument
makes their divisor principal; the remaining class is constant;
and a curve through the fixed rational point and an arbitrary
geometric point proves pointwise equality with one global unit.
These hypotheses all hold here. Finally Carlson--Haine--Wolf,
Theorem~4.18, applies to every regular affine source and all its
sliced etale objects, giving (F50). Thus the conclusion concerns
the full relative geometric morphisms after localization, including
constants and noninvertible maps, not only their point functions
or their fundamental-group homomorphisms.
\end{proof}
